\section{Independence and dependence}

Conditional probability lets us ask how the probability of one event changes
after another event is known. The next natural question is therefore:
\[
\text{when does information change the probability, and when does it not?}
\]

This is the conceptual origin of independence. Independence is not first of all
a multiplication formula. It is a statement that learning one event occurred
does not change the probability of another event.

\subsection*{Information that changes probability}

Draw one card from a standard deck. Let
\[
A=\{\text{the card is an ace}\}.
\]
Before any information is given,
\[
P(A)=\frac4{52}=\frac1{13}.
\]

Now suppose we are told that the card is a face card. Let
\[
B=\{\text{the card is a face card}\}.
\]
Inside the face cards there are no aces, using the standard convention that face
cards are jacks, queens, and kings. Therefore
\[
P(A\mid B)=0.
\]
The information \(B\) changes the probability of \(A\). It changes it from
\(1/13\) to \(0\). In this case, the events are dependent.

\subsection*{Information that does not change probability}

Now let
\[
C=\{\text{the card is a heart}\}.
\]
If we are told that the card is a heart, the relevant space consists of the
\(13\) hearts. Among them, exactly one is an ace. Hence
\[
P(A\mid C)=\frac1{13}.
\]
But this is the same as the original probability \(P(A)\). Knowing that the card
is a heart does not change the probability that it is an ace.

This is the basic meaning of independence:
\[
\text{the information occurred, but it did not change the probability of the event of interest.}
\]

\begin{definitionbox}
Events \(A\) and \(B\), with \(P(B)>0\), are \textbf{independent in the
conditional sense} if
\[
P(A\mid B)=P(A).
\]
They are \textbf{dependent} if knowing \(B\) changes the probability of \(A\),
that is, if
\[
P(A\mid B)\neq P(A).
\]
\end{definitionbox}

This definition expresses the intuitive meaning directly. It says that \(B\)
brings no probabilistic information about \(A\).

\subsection*{The multiplication form of independence}

The conditional definition is conceptually clear, but it requires \(P(B)>0\).
Using the definition of conditional probability, the equality
\[
P(A\mid B)=P(A)
\]
becomes
\[
\frac{P(A\cap B)}{P(B)}=P(A).
\]
Multiplying both sides by \(P(B)\), we obtain
\[
P(A\cap B)=P(A)P(B).
\]

This gives the standard symmetric definition.

\begin{definitionbox}
Events \(A\) and \(B\) are \textbf{independent} if
\[
P(A\cap B)=P(A)P(B).
\]
If this equality fails, the events are \textbf{dependent}.
\end{definitionbox}

This form is symmetric in \(A\) and \(B\). It makes clear that if \(A\) gives no
information about \(B\), then \(B\) gives no information about \(A\), provided
the relevant conditional probabilities are defined.

\begin{remarkbox}
The product formula should not be treated as the first meaning of independence.
It is the algebraic form of the idea that conditioning on one event does not
change the probability of the other.
\end{remarkbox}

\subsection*{Example: two coin tosses}

Toss a fair coin twice and record the ordered sequence:
\[
\Omega=\{HH,HT,TH,TT\}.
\]
Let
\[
A=\{\text{the first toss is heads}\}=\{HH,HT\},
\]
and
\[
B=\{\text{the second toss is heads}\}=\{HH,TH\}.
\]
Then
\[
P(A)=\frac12,\qquad P(B)=\frac12.
\]
The intersection is
\[
A\cap B=\{HH\},
\]
so
\[
P(A\cap B)=\frac14.
\]
Since
\[
P(A)P(B)=\frac12\cdot\frac12=\frac14,
\]
we have
\[
P(A\cap B)=P(A)P(B).
\]
Thus \(A\) and \(B\) are independent.

The conditional interpretation says the same thing. If we know the second toss
is heads, the remaining possible sequences are
\[
HH,TH.
\]
Among these, exactly one has first toss heads. Hence
\[
P(A\mid B)=\frac12=P(A).
\]
The information about the second toss does not change the probability of the
first toss.

\subsection*{Example: cards and dependence}

Draw one card from a deck. Let
\[
A=\{\text{the card is an ace}\},
\]
and
\[
B=\{\text{the card is a face card}\}.
\]
Then
\[
A\cap B=\varnothing,
\]
so
\[
P(A\cap B)=0.
\]
But
\[
P(A)P(B)=\frac4{52}\cdot\frac{12}{52}>0.
\]
Therefore
\[
P(A\cap B)\neq P(A)P(B).
\]
The events are dependent. Knowing that the card is a face card makes the event
``ace'' impossible.

\subsection*{Disjointness is not independence}

Students often confuse disjoint events with independent events. These ideas are
very different.

Events \(A\) and \(B\) are disjoint if
\[
A\cap B=\varnothing.
\]
This means they cannot occur together.

Events \(A\) and \(B\) are independent if
\[
P(A\cap B)=P(A)P(B).
\]
This means that knowing one occurred does not change the probability of the
other.

If \(A\) and \(B\) are disjoint and both have positive probability, then they
cannot be independent. Indeed,
\[
P(A\cap B)=0,
\]
but
\[
P(A)P(B)>0.
\]
So the equality required for independence fails.

\begin{remarkbox}
If two positive-probability events are disjoint, then learning that one occurred
makes the other impossible. That is the strongest possible form of dependence,
not independence.
\end{remarkbox}

\subsection*{Independence is not causation}

Dependence means that probabilities change under information. It does not by
itself identify a causal mechanism. If
\[
P(A\mid B)\neq P(A),
\]
then \(B\) is informative about \(A\), but this does not automatically mean that
\(B\) causes \(A\).

For example, suppose \(A\) is the event that a randomly selected student passes
an exam and \(B\) is the event that the student attended review sessions. If
\(P(A\mid B)>P(A)\), then attendance is associated with a higher pass
probability in the model or data. But this observation alone does not prove that
attendance caused the result. There may be other factors, such as preparation,
motivation, or prior knowledge.

Probability detects changes of likelihood under information. Causation requires
additional assumptions and is a separate question.

\subsection*{Independence of repeated trials as a modelling assumption}

When we say that repeated coin tosses are independent, we are making a modelling
assumption. For two tosses, this assumption says, for example, that knowing the
first toss was heads does not change the probability that the second toss is
heads.

If \(P(H)=p\) for each toss and the tosses are independent, then
\[
P(HH)=p^2,
\]
\[
P(HT)=p(1-p),
\]
\[
P(TH)=(1-p)p,
\]
\[
P(TT)=(1-p)^2.
\]

Without the independence assumption, we cannot automatically multiply the
single-toss probabilities. The second toss might depend on the first in some
artificial mechanism. Independence is not guaranteed by notation such as
\(\{H,T\}^2\); it must be part of the probability assignment.

\begin{examplebox}
Suppose a device produces two symbols, each either \(H\) or \(T\), and the sample
space is \(\{HH,HT,TH,TT\}\). If the device always repeats the first symbol, then
only \(HH\) and \(TT\) have positive probability. The sample space still uses two
positions, but the two positions are not independent.
\end{examplebox}

\subsection*{More than two events}

For several events, independence becomes more subtle. It is not enough that each
pair behaves independently. For three events \(A,B,C\), full mutual independence
requires the product rule for every relevant intersection:
\[
P(A\cap B)=P(A)P(B),
\]
\[
P(A\cap C)=P(A)P(C),
\]
\[
P(B\cap C)=P(B)P(C),
\]
and also
\[
P(A\cap B\cap C)=P(A)P(B)P(C).
\]

At this stage, the main point is conceptual: independence means that information
from some events does not change probabilities of the others. For many events,
this must hold consistently across combinations of events.

\subsection*{What independence teaches}

Independence is a statement about information. It says that the occurrence of
one event does not alter the probability assigned to another event. The product
formula is important because it gives an algebraic test, but the meaning remains
conditional:
\[
\text{knowing }B\text{ does not change the probability of }A.
\]

\begin{theorembox}
For events of positive probability, the following ideas are equivalent:
\[
P(A\mid B)=P(A),
\]
\[
P(B\mid A)=P(B),
\]
\[
P(A\cap B)=P(A)P(B).
\]
Each expresses the same idea: the events do not provide probabilistic
information about each other.
\end{theorembox}
